\documentclass{amsart}
% For dealing with references we use the comment environment
\usepackage{verbatim}
\newenvironment{reference}{\comment}{\endcomment}
%\newenvironment{reference}{}{}
\newenvironment{slogan}{\comment}{\endcomment}
\newenvironment{history}{\comment}{\endcomment}

% For commutative diagrams we use Xy-pic
\usepackage[all]{xy}

% We use 2cell for 2-commutative diagrams.
\xyoption{2cell}
\UseAllTwocells

% We use multicol for the list of chapters between chapters
\usepackage{multicol}

% This is generally recommended for better output
\usepackage{lmodern}
\usepackage[T1]{fontenc}

% For cross-file-references

% Package for hypertext links:
\usepackage{hyperref}

% For any local file, say "hello.tex" you want to link to please

% Theorem environments.
%
\theoremstyle{plain}
\newtheorem{theorem}[subsection]{Theorem}
\newtheorem{proposition}[subsection]{Proposition}
\newtheorem{lemma}[subsection]{Lemma}

\theoremstyle{definition}
\newtheorem{definition}[subsection]{Definition}
\newtheorem{example}[subsection]{Example}
\newtheorem{exercise}[subsection]{Exercise}
\newtheorem{situation}[subsection]{Situation}

\theoremstyle{remark}
\newtheorem{remark}[subsection]{Remark}
\newtheorem{remarks}[subsection]{Remarks}

\numberwithin{equation}{subsection}

% Macros
%
\def\lim{\mathop{\mathrm{lim}}\nolimits}
\def\colim{\mathop{\mathrm{colim}}\nolimits}
\def\Spec{\mathop{\mathrm{Spec}}}
\def\Hom{\mathop{\mathrm{Hom}}\nolimits}
\def\Ext{\mathop{\mathrm{Ext}}\nolimits}
\def\SheafHom{\mathop{\mathcal{H}\!\mathit{om}}\nolimits}
\def\SheafExt{\mathop{\mathcal{E}\!\mathit{xt}}\nolimits}
\def\Sch{\mathit{Sch}}
\def\Mor{\mathop{\mathrm{Mor}}\nolimits}
\def\Ob{\mathop{\mathrm{Ob}}\nolimits}
\def\Sh{\mathop{\mathit{Sh}}\nolimits}
\def\NL{\mathop{N\!L}\nolimits}
\def\CH{\mathop{\mathrm{CH}}\nolimits}
\def\proetale{{pro\text{-}\acute{e}tale}}
\def\etale{{\acute{e}tale}}
\def\QCoh{\mathit{QCoh}}
\def\Ker{\mathop{\mathrm{Ker}}}
\def\Im{\mathop{\mathrm{Im}}}
\def\Coker{\mathop{\mathrm{Coker}}}
\def\Coim{\mathop{\mathrm{Coim}}}

% Boxtimes
%
\DeclareMathSymbol{\boxtimes}{\mathbin}{AMSa}{"02}

%
% Macros for moduli stacks/spaces
%
\def\QCohstack{\mathcal{QC}\!\mathit{oh}}
\def\Cohstack{\mathcal{C}\!\mathit{oh}}
\def\Spacesstack{\mathcal{S}\!\mathit{paces}}
\def\Quotfunctor{\mathrm{Quot}}
\def\Hilbfunctor{\mathrm{Hilb}}
\def\Curvesstack{\mathcal{C}\!\mathit{urves}}
\def\Polarizedstack{\mathcal{P}\!\mathit{olarized}}
\def\Complexesstack{\mathcal{C}\!\mathit{omplexes}}
% \Pic is the operator that assigns to X its picard group, usage \Pic(X)
% \Picardstack_{X/B} denotes the Picard stack of X over B
% \Picardfunctor_{X/B} denotes the Picard functor of X over B
\def\Pic{\mathop{\mathrm{Pic}}\nolimits}
\def\Picardstack{\mathcal{P}\!\mathit{ic}}
\def\Picardfunctor{\mathrm{Pic}}
\def\Deformationcategory{\mathcal{D}\!\mathit{ef}}

\setlength{\emergencystretch}{3em}
\begin{document}
\title[Stacks Project, Tag 0CT9]{484 --- Proper open-and-closed fibre pieces lift as proper open-and-closed pieces over henselian pairs in separated finite-type families}
\maketitle
\noindent Copyright (C) 2005--2025 Johan de Jong. Stacks Project authors. Source: \href{https://stacks.math.columbia.edu/tag/0CT9}{Tag 0CT9}.

\noindent Editorial difficulty note (not author text). Difficulty H2: Chow modification and projective closure must transport the desired component without introducing an overlapping boundary. Henselianity rules out the closed intersections and boundary pieces, after which Stein factorization and lifting of integral-algebra idempotents construct the proper component..

\noindent Editorial provenance note (not author text). History: excerpt from revision \texttt{a0444\allowbreak 6e57e\allowbreak c1fbc\allowbreak 252a8\allowbreak 71afc\allowbreak ec775\allowbreak 2fb28\allowbreak 07b14}, more-morphisms.tex, lines 15462--15547. Reference presentation was adapted; original-author context and core support blocks were retained or added. The mathematical wording is unchanged. Licensed under GNU FDL 1.2 or later; no invariant sections or cover texts. See ../licenses/COPYING. Context blocks reproduce author definitions and supporting results; they are not additional counted problems.

\begin{lemma}
\label{lemma-split-off-proper-part-henselian}
\begin{reference}
A reference for the case of an adic Noetherian base is
\href{https://stacks.math.columbia.edu/bibliography}{Bibliography: EGA (III, Proposition 5.5.1)}
\end{reference}
Let $(A, I)$ be a henselian pair. Let $X \to \Spec(A)$
be separated and of finite type. Set $X_0 = X \times_{\Spec(A)} \Spec(A/I)$.
Let $Y \subset X_0$ be an open and closed subscheme such that
$Y \to \Spec(A/I)$ is proper. Then there exists an open and closed
subscheme $W \subset X$ which is proper over $A$ with
$W \times_{\Spec(A)} \Spec(A/I) = Y$.
\end{lemma}

\begin{proof}
We will denote $T \mapsto T_0$ the base change by $\Spec(A/I) \to \Spec(A)$.
By Chow's lemma (in the form of
Limits, Lemma \href{https://stacks.math.columbia.edu/tag/0202}{lemma chow finite type (Tag 0202)})
there exists a surjective proper morphism $\varphi : X' \to X$ such
that $X'$ admits an immersion into $\mathbf{P}^n_A$.
Set $Y' = \varphi^{-1}(Y)$. This is an open and closed subscheme
of $X'_0$. Suppose the lemma holds for $(X', Y')$. Let $W' \subset X'$
be the open and closed subscheme proper over $A$ such that $Y' = W'_0$.
By Morphisms, Lemma \href{https://stacks.math.columbia.edu/tag/01W6}{lemma image proper scheme closed (Tag 01W6)}
$W = \varphi(W') \subset X$ and
$Q = \varphi(X' \setminus W') \subset X$ are closed subsets and by
Morphisms, Lemma \href{https://stacks.math.columbia.edu/tag/03GN}{lemma image proper is proper (Tag 03GN)}
$W$ is proper over $A$. The image of $W \cap Q$ in $\Spec(A)$ is closed.
Since $(A, I)$ is henselian, if $W \cap Q$ is nonempty, then we
find that $W \cap Q$ has a point lying over $\Spec(A/I)$.
This is impossible as $W'_0 = Y' = \varphi^{-1}(Y)$.
We conclude that $W$ is an open and closed subscheme
of $X$ proper over $A$ with $W_0 = Y$.
Thus we reduce to the case described in the next paragraph.

\medskip\noindent
Assume there exists an immersion $j : X \to \mathbf{P}^n_A$ over $A$.
Let $\overline{X}$ be the scheme theoretic image of $j$.
Since $j$ is a quasi-compact morphism
(Schemes, Lemma \href{https://stacks.math.columbia.edu/tag/03GI}{lemma quasi compact permanence (Tag 03GI)})
we see that $j : X \to \overline{X}$ is an open immersion
(Morphisms, Lemma \href{https://stacks.math.columbia.edu/tag/01RG}{lemma quasi compact immersion (Tag 01RG)}).
Hence the base change $j_0 : X_0 \to \overline{X}_0$
is an open immersion as well.
Thus $j_0(Y) \subset \overline{X}_0$ is open.
It is also closed by Morphisms, Lemma
\href{https://stacks.math.columbia.edu/tag/01W6}{lemma image proper scheme closed (Tag 01W6)}.
Suppose that the lemma holds for $(\overline{X}, j_0(Y))$.
Let $\overline{W} \subset \overline{X}$ be the
corresponding open and closed subscheme proper over $A$
such that $j_0(Y) = \overline{W}_0$.
Then $T = \overline{W} \setminus j(X)$ is closed in $\overline{W}$,
hence has closed image in $\Spec(A)$ by properness of $\overline{W}$
over $A$. Since $(A, I)$ is henselian, we find that if $T$
is nonempty, then there is a point of $T$ mapping into $\Spec(A/I)$.
This is impossible because $j_0(Y) = \overline{W}_0$ is contained in $j(X)$.
Hence $\overline{W}$ is contained in $j(X)$ and we can
set $W \subset X$ equal to the unique open and closed
subscheme mapping isomorphically to $\overline{W}$ via $j$.
Thus we reduce to the case described in the next paragraph.

\medskip\noindent
Assume $X \subset \mathbf{P}^n_A$ is a closed subscheme.
Then $X \to \Spec(A)$ is a proper morphism.
Let $Z = X_0 \setminus Y$. This is an open and closed
subscheme of $X_0$ and $X_0 = Y \amalg Z$.
Let $X \to X' \to \Spec(A)$ be the Stein factorization as in
Theorem \href{https://stacks.math.columbia.edu/tag/03H2}{theorem stein factorization general (Tag 03H2)}.
Let $Y' \subset X'_0$ and $Z' \subset X'_0$ be the images of
$Y$ and $Z$.
Since the fibres of $X \to Z$ are geometrically connected,
we see that $Y' \cap Z' = \emptyset$.
Hence $X'_0 = Y' \amalg Z'$ as $X \to X'$ is surjective.
Since $X' \to \Spec(A)$ is integral, we see that
$X'$ is the spectrum of an $A$-algebra integral over $A$.
Recall that open and closed subsets of spectra correspond
$1$-to-$1$ with idempotents in the corresponding ring, see
Algebra, Lemma \href{https://stacks.math.columbia.edu/tag/00EE}{lemma disjoint decomposition (Tag 00EE)}.
Hence by
More on Algebra, Lemma \href{https://stacks.math.columbia.edu/tag/09XI}{lemma characterize henselian pair (Tag 09XI)}
we see that we may write $X' = W' \amalg V'$
with $W'$ and $V'$ open and closed and
with $Y' = W'_0$ and $Z' = V'_0$.
Let $W$ be the inverse image in $X$
to finish the proof.
\end{proof}
\end{document}
